\chapter{System Design and Architecture}
\label{ch:system}

\section{Overview}

The system is designed as a multi-layered pipeline in which each layer takes the output of
the previous layer as input, progressively transforming a raw label image into a structured
validation report. \Cref{fig:pipeline_overview} illustrates the top-level architecture.

\begin{figure}[h]
\centering
\begin{tikzpicture}[
  node distance=0.8cm,
  box/.style={rectangle, rounded corners=3pt, draw=black!60, fill=blue!8,
              text width=3.8cm, minimum height=1cm, align=center, font=\small},
  arrow/.style={-Stealth, thick},
  label/.style={font=\scriptsize, text=gray}
]

\node[box] (input)   {Input\\(PDF / Image)};
\node[box, right=of input]  (extract) {Extraction Layer\\OCR + Barcodes};
\node[box, right=of extract] (type)   {Type Classifier\\(schema scoring)};
\node[box, right=of type]   (valid)   {Validator\\(rule engine)};
\node[box, right=of valid]  (report)  {Report\\(JSON + UI)};

\draw[arrow] (input)   -- (extract);
\draw[arrow] (extract) -- (type);
\draw[arrow] (type)    -- (valid);
\draw[arrow] (valid)   -- (report);

\node[label, below=0.3cm of input]   {raw bytes};
\node[label, below=0.3cm of extract] {ExtractedLabel};
\node[label, below=0.3cm of type]    {schema match};
\node[label, below=0.3cm of valid]   {RuleResult list};

\end{tikzpicture}
\caption{Top-level pipeline architecture. Each stage produces a typed data structure
consumed by the next stage.}
\label{fig:pipeline_overview}
\end{figure}

The pipeline is exposed via a FastAPI \ac{REST} \ac{API} (\texttt{/api/extract},
\texttt{/api/merge-validate}) and via a Streamlit \ac{UI} that allows QA operators to
upload a label, inspect the extracted fields, apply manual corrections, and trigger
validation interactively.

\section{Extraction Layer}

The extraction layer is responsible for converting a raw label—provided as a PDF or raster
image (PNG/JPG)—into an \texttt{ExtractedLabel} data structure that holds all fields the
validator requires.

\subsection{Ingestion and Preprocessing}

The ingestion module (\texttt{LabelExtractor}) accepts a file path and DPI parameter.
PDF inputs are rendered to high-resolution raster images using pdf2image (poppler backend)
at the specified DPI (default~300). Raster inputs are passed directly to the preprocessing
stage.

Preprocessing applies:
\begin{enumerate}
  \item Contrast enhancement via CLAHE (Contrast Limited Adaptive Histogram Equalization).
  \item Adaptive Gaussian thresholding to binarise text regions without losing barcode gradient
  information.
  \item Image sharpening via an unsharp mask kernel.
\end{enumerate}

\subsection{Zone Segmentation}

The label is divided into named zones by anchoring on detected barcode positions. A
DataMatrix barcode, when detected, is used to calibrate the centre zone; the identity zone
is estimated to occupy the left portion, and date/compliance zones occupy the right portion.
Zone boundaries are stored as normalised fractions of image height and width, allowing the
same zone definitions to apply regardless of image resolution.

\subsection{PDF-Native Text Extraction (Pass 1)}

For PDF inputs, PyMuPDF (fitz) is used to extract text spans with their exact page
coordinates, font names, font sizes, and character-level detail. Each span is converted to
an \texttt{OcrWord}-equivalent data structure with confidence~1.0. This pass produces
near-perfect field extraction on programmatically-generated label PDFs, as the text is
embedded as vector data.

\subsection{OCR Pass (Pass 5)}

For raster inputs—or for regions of a PDF where text is embedded as images—the \ac{OCR}
ensemble is invoked. Tesseract~5.3 and EasyOCR~1.7 are run in sequence on the
preprocessed image. The resulting word lists are merged by late fusion:

\begin{enumerate}
  \item Bounding boxes from both engines are compared pairwise.
  \item If two boxes overlap with IoU $>$ 0.4, they are considered candidates for the same
  word.
  \item The candidate with higher confidence is retained; if confidence values are within
  0.05 of each other, the Tesseract result is preferred (it tends to produce cleaner
  article numbers and date strings on IKEA label imagery).
  \item Non-overlapping boxes from both engines are included without modification.
\end{enumerate}

Words with confidence below~0.25 are filtered out before field mapping to avoid
propagating garbage tokens to the rule engine.

\subsection{Barcode Decoding}

Three barcode decoding calls are made in sequence:

\begin{enumerate}
  \item \textbf{pyzbar} for ITF-14 and EAN-13: applied to a grayscale copy of the
  preprocessed image. A morphological dilation pass is applied to thin-bar ITF codes
  at 25\% magnification, which would otherwise fall below pyzbar's minimum module width.
  \item \textbf{pylibdmtx} for DataMatrix: applied to a region crop around the centre
  zone. If the first attempt fails, a rotation sweep (0°, 90°, 180°, 270°) is tried.
  \item GS1 parsing of the decoded DataMatrix payload by the custom \texttt{gs1\_parser}
  module, which populates the AI-indexed fields on the \texttt{ExtractedLabel}.
\end{enumerate}

\section{Field Mapper}

The field mapper (\texttt{field\_mapper.py}) converts the merged word list and barcode
results into a fully populated \texttt{ExtractedLabel}. It applies a cascade of named
regular-expression patterns to the sorted word stream, assigning each matched token to the
corresponding field. The key patterns are listed in \cref{tab:regex_fields}.

\begin{table}[h]
\centering
\caption{Selected regex patterns used in field extraction.}
\label{tab:regex_fields}
\begin{tabularx}{\textwidth}{lX}
\toprule
\textbf{Field} & \textbf{Pattern} \\
\midrule
Article number & \texttt{\textbackslash b(\textbackslash d\{3\})[\textbackslash s.\textbackslash-](\textbackslash d\{3\})[\textbackslash s.\textbackslash-](\textbackslash d\{2\})\textbackslash b} \\
Plant ID & \texttt{\textbackslash bPI[-–]\textbackslash d\{3,9\}[-–]\textbackslash d\{1,3\}\textbackslash b} \\
Human date & \texttt{(\textbackslash d\{2\}[-./]\textbackslash d\{2\}[-./]\textbackslash d\{2\})...(YY-?MM-?DD)} \\
Weight (kg) & \texttt{\textbackslash b(\textbackslash d+(?:\textbackslash.\textbackslash d+)?)\textbackslash s*k(?:g)?\textbackslash b} \\
Supplier address & \texttt{SE\textbackslash s*[-–]\textbackslash s*\textbackslash d\{5\}\textbackslash s*\textbackslash S+} \\
Copyright & \texttt{(?:©|\textbackslash(c\textbackslash)|...)?Inter\textbackslash s+IKEA[\textbackslash s\textbackslash S]\{0,180\}?B\textbackslash.?V\textbackslash.?} \\
\bottomrule
\end{tabularx}
\end{table}

\section{Label-Type Classifier}

\subsection{Motivation}

The validation rules differ between label types—a rule that is mandatory for \texttt{1R35-IDADM\_IDBDM}
may be irrelevant for \texttt{1CIR24-PRADM}. The classifier must therefore identify the
correct schema before the validator can run. With 34 schemas in the corpus, a brute-force
approach of validating against all schemas and choosing the best-scoring one would generate
misleading multi-schema results.

\subsection{Scoring Heuristic}

The classifier computes a weighted field-presence score for each schema candidate:

\begin{equation}
S(l, s) = \frac{\displaystyle\sum_{f \in R_s} \mathbf{1}[\text{present}(l, f)]
              + 0.4 \cdot \sum_{f \in O_s} \mathbf{1}[\text{present}(l, f)]
              + B(l, s)}
           {\max(|R_s| + 0.4|O_s| + 3,\; 1)}
\label{eq:type_score}
\end{equation}

where $R_s$ and $O_s$ are the required and optional field lists of schema $s$,
$\mathbf{1}[\text{present}(l, f)]$ is~1 if field $f$ is present in the extraction $l$,
and $B(l, s)$ is a barcode bonus/penalty term:

\begin{equation}
B(l, s) = \sum_{b \in \{\text{ITF}, \text{DM}, \text{EAN}\}} \left(
  \mathbf{1}[\text{schema has } b] \cdot \mathbf{1}[\text{label has } b]
  - 0.5 \cdot \mathbf{1}[\text{schema has } b] \cdot \mathbf{1}[\text{label lacks } b]
\right)
\label{eq:barcode_bonus}
\end{equation}

The schema with the highest normalised score $S \in [0, 1]$ is selected as the best match.
The top-5 shortlist is also returned to the \ac{UI} so a QA operator can override the
automatic selection if needed.

\section{Rule-Based Validator}

\subsection{Rule Types}

The validator supports four rule check types:

\begin{description}
  \item[\texttt{present}] The specified field must be extracted (non-None, non-empty
  after stripping whitespace). Failure indicates a missing required element.

  \item[\texttt{regex\_match}] The field value must fully match a provided regular
  expression pattern. Used for structured fields such as article numbers (\texttt{\textbackslash d\{3\}.\textbackslash d\{3\}.\textbackslash d\{2\}}), origin texts, and date stamps.

  \item[\texttt{exact\_match}] The field value must match an expected string exactly
  (accent-normalised, case-folded). Used for supplier name and copyright notice fields
  that are the same on every label.

  \item[\texttt{ai\_field\_present}] A specific GS1 Application Identifier must be present
  in the decoded DataMatrix payload. Failure can indicate either that the DataMatrix was not
  decoded at all or that it was decoded but does not contain the required AI.
\end{description}

\subsection{Compliance Scoring}

Each label receives a weighted compliance score:

\begin{equation}
\text{CompScore}(l, s) = \frac{\displaystyle\sum_{r \in \text{runnable}} w_r \cdot \mathbf{1}[\text{pass}(r)]}
           {\displaystyle\sum_{r \in \text{runnable}} w_r}
\label{eq:compliance_score}
\end{equation}

where $w_r = 2.0$ for error-severity rules and $w_r = 1.0$ for warning-severity rules,
and ``runnable'' excludes skipped rules (e.g.\ physical-dimension checks). A
CompScore of 1.0 means all runnable rules passed; a score below 0.5 is considered a
definitive non-compliance recommendation.

\subsection{Additional Checks}

Beyond field-level rules, the validator architecture supports:
\begin{itemize}
  \item \textbf{Layout check} (\texttt{layout\_checker.py}): spatial zone compliance.
  Currently returns an empty report (stub); intended as future work.
  \item \textbf{Overlap detection} (\texttt{overlap\_detector.py}): identifies extracted
  field bounding boxes that unexpectedly overlap.
  \item \textbf{Completeness check} (\texttt{completeness.py}): cross-field consistency
  validation (DataMatrix date vs.\ human-readable date). Currently a stub.
\end{itemize}

\section{REST API}

The pipeline is exposed via a FastAPI application with three primary endpoints:

\begin{description}
  \item[\texttt{POST /api/extract}] Accepts a label file (PDF or image) and DPI parameter.
  Returns the \texttt{ExtractedLabel} as a JSON object, a SHA-256 file hash for
  deduplication, and a base-64-encoded preview PNG.

  \item[\texttt{POST /api/merge-validate}] Accepts an \texttt{ExtractedLabel} JSON and
  an optional corrections dictionary (from the UI's human-correction step). Applies the
  corrections, runs the full validator, and returns a \texttt{ValidationReport}.

  \item[\texttt{POST /api/export-jsonl}] Exports the extracted+corrected label as an
  NDJSON training record for future fine-tuning of field extraction models.
\end{description}

All endpoints accept CORS from any origin, allowing the React/TypeScript frontend to
communicate with the local API during development.

\section{Interactive UI}

The Streamlit \ac{UI} (\texttt{app.py}) provides:
\begin{itemize}
  \item A drag-and-drop file uploader for PDF and image inputs.
  \item A field-by-field display of extracted values with editable correction inputs.
  \item A colour-coded validation report showing passed rules (green), warnings (amber),
  and errors (red).
  \item A type-confidence bar chart showing the top-5 schema candidates and their scores.
  \item A canvas-based bounding-box overlay highlighting where each field was detected.
\end{itemize}

The React/TypeScript frontend (\texttt{frontend/}) provides an alternative \ac{UI}
for tighter integration with IKEA's existing tooling, served via the FastAPI backend.
